\documentclass[11pt,a4paper]{article}

\usepackage[T1]{fontenc}
\usepackage[utf8]{inputenc}
\usepackage[margin=1in]{geometry}
\usepackage{amsmath}
\usepackage{amssymb}
\usepackage{booktabs}
\usepackage{microtype}
\usepackage[colorlinks=true,linkcolor=blue,urlcolor=blue]{hyperref}

\newcommand{\nnz}{\operatorname{nnz}}
\newcommand{\epsm}{\varepsilon_{\mathrm{machine}}}

\title{Complexity of Eigensolvers versus Linear Solvers\\
       \large Dense and Sparse Cases Based on Trefethen and Bau}
\author{}
\date{}

\begin{document}

\maketitle

\begin{abstract}
This report compares the arithmetic complexity of linear-system solvers and
eigenvalue solvers using the discussion and operation counts in Trefethen and
Bau's \emph{Numerical Linear Algebra}.  For dense matrices, a direct linear
solve and a complete eigenvalue computation are both cubic in the matrix order,
although their leading constants and outputs differ.  For sparse matrices,
Krylov linear solvers and partial eigensolvers have a similar
matrix--vector-product-dominated complexity structure.  This similarity is
conditional: it does not imply equal iteration counts, equal elapsed time, or a
universal common complexity exponent.
\end{abstract}

\noindent
All citations below refer to the printed pagination of the book.  In the local
scan
\path|docs/Bau_Numerical_linear_Algebra_BOOK.pdf|, the PDF page number is the
printed page number plus~10.  The matrix order is denoted by $m$.

\section{Dense direct algorithms}

For one dense general system $Ax=b$, Gaussian elimination computes an LU
factorization in approximately
\begin{equation}
  \frac{2}{3}m^3
\end{equation}
flops.  Forward and backward substitution then require approximately $m^2$
flops each.  Thus the leading cost of one solve is $2m^3/3$; subsequent
right-hand sides cost approximately $2m^2$ each once the factors have been
formed \cite[pp.~151--152, especially (20.8)]{trefethenbau}.  For a Hermitian
positive-definite matrix, Cholesky factorization costs approximately $m^3/3$
flops \cite[p.~175]{trefethenbau}.

The book describes a general-purpose dense eigensolver as a two-phase process.
A general matrix is first reduced to upper-Hessenberg form, and an iterative QR
phase then reduces the Hessenberg matrix to Schur form.  The first phase and the
complete computation both require $O(m^3)$ work
\cite[pp.~193--194]{trefethenbau}.  More precisely, Hessenberg reduction costs
approximately $10m^3/3$ flops \cite[p.~199, (26.1)]{trefethenbau}.

For a Hermitian matrix, Phase~1 is a reduction to tridiagonal form costing
approximately
\begin{equation}
  \frac{4}{3}m^3
\end{equation}
flops \cite[p.~200, (26.2)]{trefethenbau}.  If only eigenvalues are required,
the structured QR phase costs $O(m^2)$ and is therefore lower order
\cite[p.~194]{trefethenbau}.  For a real symmetric matrix, the resulting
eigenvalue-only algorithm is quoted as requiring approximately $4m^3/3$ flops
\cite[p.~223]{trefethenbau}.

The principal dense comparison is consequently
\begin{equation}
  \boxed{\text{dense direct linear solve}=\Theta(m^3),\qquad
         \text{dense complete eigensolve}=\Theta(m^3).}
\end{equation}
The exponents agree, but the leading constants do not.  Comparing the book's
counts, a real symmetric eigenvalue-only computation is approximately twice
the leading cost of an LU solve and four times the factorization cost of a
Cholesky solve.  These ratios are deductions from the quoted operation counts,
not separate claims made by the authors.

Computing eigenvectors changes the comparison materially.  For a real
symmetric problem, the book gives approximately $9m^3$ flops for the QR
algorithm when eigenvectors are accumulated.  Divide-and-conquer reduces this
to approximately $4m^3$, and numerical deflation often reduces the observed
combined count to about $8m^3/3$ \cite[p.~232]{trefethenbau}.  A statement about
the cost of ``an eigensolve'' must therefore specify whether only eigenvalues or
also eigenvectors are required.

Although an eigenvalue algorithm is iterative in principle, its dependence on
the target machine precision is weak.  The book notes a dependence as weak as
$\log\lvert\log\epsm\rvert$ for superlinear convergence and states that, in
practice, eigenvalue computation differs from linear-system solution by a
small constant factor, typically closer to one than to ten
\cite[p.~193 and Exercise~25.2]{trefethenbau}.

\section{Sparse matrices and matrix--vector products}

Let
\begin{equation}
  s=\nnz(A)
\end{equation}
be the number of stored nonzero entries.  If a sparse matrix has an average of
$\nu$ nonzeros per row, then $s\simeq \nu m$.  The book states that a sparse
matrix--vector product can be evaluated in
\begin{equation}
  O(\nu m)=O(s)
\end{equation}
operations instead of the $O(m^2)$ operations required for an unstructured
dense matrix \cite[pp.~244--245]{trefethenbau}.  Krylov methods exploit a matrix
through precisely this $x\mapsto Ax$ ``black box.''

The same section places linear and eigenvalue algorithms in parallel:
conjugate gradients (CG) and Lanczos apply to Hermitian problems, whereas
GMRES and Arnoldi apply to non-Hermitian problems
\cite[pp.~245--246]{trefethenbau}.  This common Krylov framework is the basis
for saying that sparse iterative linear solvers and sparse partial eigensolvers
have a similar complexity structure.

\subsection{Meaning of the iteration counts}

Define
\begin{itemize}
  \item $k_{\mathrm{lin}}$ as the number of iterations required by an
        iterative linear solver, such as CG or GMRES, to solve $Ax=b$ to a
        prescribed tolerance;
  \item $k_{\mathrm{eig}}$ as the number of iterations required by an
        iterative eigensolver, such as Lanczos or Arnoldi, to approximate the
        requested eigenpair or eigenpairs to a prescribed tolerance.
\end{itemize}

These counts are not generally known in advance.  The book explains that CG
can converge rapidly when the eigenvalues are favorably clustered away from
the origin, whereas Lanczos can find selected eigenvalues rapidly when they are
well separated from the remainder of the spectrum.  Both convergence questions
are connected with polynomial approximation, and both may be affected strongly
by preconditioning \cite[p.~246]{trefethenbau}.

\subsection{Sparse iterative linear solve}

A CG iteration contains one multiplication by $A$ and several vector
operations.  The book gives approximately $2m^2$ flops per iteration for a
dense unstructured matrix, but only $O(m)$ per iteration when sparsity or other
structure makes the matrix--vector product linear-time
\cite[p.~295]{trefethenbau}.  In terms of the number of nonzeros, the leading
model is
\begin{equation}
  C_{\mathrm{linear}}=O\!\left(k_{\mathrm{lin}}s\right).
  \label{eq:linear-sparse}
\end{equation}

\subsection{Sparse partial eigensolve}

Each Lanczos step consists of a matrix--vector multiplication, an inner
product, and a small number of vector operations.  The book emphasizes that
when sparsity makes matrix--vector multiplication cheap, Lanczos can be used
for matrices with dimensions in the tens or hundreds of thousands
\cite[p.~278]{trefethenbau}.  For a small number of desired eigenpairs, the
leading model is therefore
\begin{equation}
  C_{\mathrm{eigen}}=O\!\left(k_{\mathrm{eig}}s\right),
  \label{eq:eigen-sparse}
\end{equation}
apart from the relatively small projected tridiagonal eigenproblems and any
reorthogonalization that is required.

For nonsymmetric matrices, Arnoldi also requires sparse matrix--vector
products, but at iteration $j$ it orthogonalizes the new vector against the
previous $j$ basis vectors; see Algorithm~33.1
\cite[pp.~252--253]{trefethenbau}.  The algorithm therefore implies the model
\begin{equation}
  C_{\mathrm{Arnoldi}}
  =O\!\left(k_{\mathrm{eig}}s+k_{\mathrm{eig}}^2m\right),
  \qquad
  \text{storage}=O(k_{\mathrm{eig}}m).
  \label{eq:arnoldi}
\end{equation}
The $k_{\mathrm{eig}}^2m$ and storage terms in
\eqref{eq:arnoldi} are deductions from the displayed Arnoldi algorithm, not
operation-count formulas explicitly printed in the book.  Restarting changes
these costs and may also change convergence.

\section{Do sparse linear and eigen solvers have the same exponent?}

Equations~\eqref{eq:linear-sparse} and \eqref{eq:eigen-sparse} show that the
two methods have the same \emph{per-iteration} dependence on sparsity.  They do
not automatically have the same total exponent because their iteration counts
may scale differently with $m$.

For example, suppose that the number of nonzeros satisfies
\begin{equation}
  s=O(m^\alpha),
\end{equation}
and the iteration counts satisfy
\begin{equation}
  k_{\mathrm{lin}}=O(m^{\beta_{\mathrm{lin}}}),
  \qquad
  k_{\mathrm{eig}}=O(m^{\beta_{\mathrm{eig}}}).
\end{equation}
Ignoring Arnoldi orthogonalization, the resulting complexities are
\begin{align}
  C_{\mathrm{linear}}&=O(m^{\alpha+\beta_{\mathrm{lin}}}),\\
  C_{\mathrm{eigen}}&=O(m^{\alpha+\beta_{\mathrm{eig}}}).
\end{align}
The exponents agree only if
$\beta_{\mathrm{lin}}=\beta_{\mathrm{eig}}$.

For the common case of $O(1)$ nonzeros per row, $s=O(m)$ and $\alpha=1$.
If both iteration counts remain bounded independently of $m$, then
\begin{equation}
  C_{\mathrm{linear}}=O(m),
  \qquad
  C_{\mathrm{eigen}}=O(m).
\end{equation}
In this restricted situation both solvers have exponent one.  If either
iteration count grows with the problem size, its growth must be included in
the exponent.

The claim also depends on comparing like outputs.  Equations
\eqref{eq:linear-sparse} and \eqref{eq:eigen-sparse} compare one linear-system
solution with a \emph{small number} of selected eigenpairs.  Computing the
complete spectrum, and especially all eigenvectors, is a different task and
cannot generally be assigned the partial Krylov cost $O(k_{\mathrm{eig}}s)$.

\section{Equal asymptotic complexity does not imply equal runtime}

Even when both methods have the same exponent, their elapsed times need not be
similar.  A more informative model is
\begin{align}
  T_{\mathrm{linear}}
    &\simeq c_{\mathrm{linear}}k_{\mathrm{lin}}s,\\
  T_{\mathrm{eigen}}
    &\simeq c_{\mathrm{eigen}}k_{\mathrm{eig}}s
      +T_{\mathrm{orth}}+T_{\mathrm{projected}},
\end{align}
where the constants and additional terms include such effects as
\begin{itemize}
  \item different iteration counts and convergence tolerances;
  \item the number of requested eigenpairs;
  \item orthogonalization, reorthogonalization, restarts, and convergence
        checks in an eigensolver;
  \item the setup and application costs of a preconditioner;
  \item memory traffic, data layout, parallelism, and implementation quality;
  \item factorization costs when shift-and-invert is used.
\end{itemize}

For example, $20m$ and $200m$ are both $O(m)$, but the second expression is ten
times larger for every $m$.  Big-$O$ notation describes growth with problem
size; it suppresses precisely the constants and lower-order terms that often
control observed runtime.

\section{Sparse direct methods and fill-in}

Sparsity does not automatically reduce a direct factorization to $O(s)$ work.
The book warns that Gaussian and Householder transformations manipulate matrix
entries and generally destroy sparsity.  It mentions nested dissection and
minimum-degree reordering as sparse direct techniques but explicitly does not
develop them \cite[p.~245]{trefethenbau}.  It also observes that Gaussian
elimination and Cholesky factorization introduce fill-in, so a factor of a
sparse matrix will usually not remain very sparse
\cite[p.~316]{trefethenbau}.

The book consequently does not provide one universal exponent for sparse
direct LU, sparse Cholesky, or a sparse direct eigensolver.  Such an exponent
depends on the sparsity graph, ordering, fill-in, and problem dimension.
Incomplete Cholesky and incomplete LU are presented as ways to retain selected
sparsity and construct useful preconditioners rather than exact factorizations
\cite[p.~316]{trefethenbau}.

Shift-and-invert deserves special mention.  In this approach an eigensolver
repeatedly applies an operator such as $(A-\sigma I)^{-1}$, so an inner linear
solve or factor application occurs inside each eigenvalue iteration.  Its cost
cannot be represented by a bare sparse matrix--vector product unless the cost
of that inner solve is included.  This observation is consistent with the
book's comparison of power, inverse, and Rayleigh-quotient iterations
\cite[p.~209]{trefethenbau}.

\section{Conclusion}

The broad unqualified statement
\begin{quote}
``For sparse matrices, linear solvers and eigensolvers have the same
asymptotic complexity''
\end{quote}
is not supported by the book because it omits the algorithm, required output,
iteration-count scaling, orthogonalization, preconditioning, and fill-in.

A defensible book-based statement is:
\begin{quote}
For sparse matrices, Krylov iterative methods for solving one linear system
and for computing a small number of eigenpairs often have a similar
asymptotic complexity structure: both are dominated by a problem-dependent
number of sparse matrix--vector products.
\end{quote}
Equivalently,
\begin{equation}
  \boxed{
  C_{\mathrm{linear}}=O(k_{\mathrm{lin}}\nnz(A)),\qquad
  C_{\mathrm{eigen}}=O(k_{\mathrm{eig}}\nnz(A))
  }
\end{equation}
for CG/Lanczos-type comparisons and a small number of eigenpairs.  Thus the
per-iteration exponents are similar; the total exponents and computational
times are not necessarily similar.

\begin{thebibliography}{1}

\bibitem{trefethenbau}
Lloyd N. Trefethen and David Bau III,
\newblock \emph{Numerical Linear Algebra},
\newblock Society for Industrial and Applied Mathematics, 1997.

\end{thebibliography}

\end{document}
